\section{Introduction}
\label{sec:intro}

Understanding the internal representations of the human mind is central to uncovering the mechanisms of cognition, from what sensory information is selectively attended to construct memory~\cite{zhan2019modelling}, to what specific information is predicted to guide adaptive decision-making~\cite{friston2009predictive,clark2013whatever,yan2023network}. Although existing models of prediction, sensory coding, memory, and decision making are implicitly framed within the information processing paradigm~\cite{marr2010vision}, precisely tracking how this processing unfolds within the 'black box' of the mind remains a considerable challenge. 

This challenge stems from the inherently subjective nature of cognition~\cite{boogert2018measuring}, which introduces both variability and ambiguity into the modeling process. First, individuals exhibit subjective experiences and idiosyncratic biases that shape how information is prioritized~\cite{zhan2021modeling}, resulting in considerable variability in internal cognitive processes~\cite{stanovich2008individual}. Second, this variability creates an ill-posed mapping between high-dimensional sensory inputs and low-dimensional behavioral outputs (e.g., identity judgments such as 'Tom'). Multiple, mechanistically distinct algorithms can yield indistinguishable behavioral responses (e.g., one observer may rely on nasal prominence and chin morphology, while another integrates nasal features with skin tone). Third, this ill-posed nature challenges conventional modeling approaches that fit canonical algorithms to population-averaged behaviors. Inferring algorithmic solutions \textit{a priori} and evaluating models solely based on their ability to simulate average behavior risks overlooking optimal, individualized solutions, thereby hindering the discovery of novel cognitive phenomena and the advancement of mechanistic theory.

 \textbf{Contributions}. We propose a data-driven framework that explicitly models trial-resolved signatures of internal information processing at the individual level. Our main contributions are:
(1) \textbf{A modular inference framework} that approximates internal information processing as Bayesian inference, with separate modules for cognitive constraints that impose bounded rationality and personalized priors reflecting heterogeneous experience[~\cite{cohen2013high,jaworska2022different}].
(2) \textbf{Trial-resolved inference} in a demanding context––novel category induction over a high-dimensional feature space––demonstrating the model’s capacity to recover subject-specific feature exploration over the course of learning. Crucially, our model operates over a theoretically unconstrained hypothesis space and produces high-fidelity, trial-wise estimates of internal processing dynamics.
(3) \textbf{Epistemological insights into computational modeling} via ablation studies, which reveal a key dissociation--models can reproduce observable behavior to a high degree while mischaracterizing internal representations. This highlights a fundamental epistemological risk of using behavioral fit alone as a proxy for cognitive validity.
(4) \textbf{Mechanistic insights into human learning behavior} through computational dissection, which identify an adaptive analysis-by-synthesis strategy: agents dynamically tune their prior-sampling distributions to balance exploratory hypothesis search with exploitative evidence gathering, offsetting memory capacity and enabling robust learning from sparse data.


 \textbf{Related works}. Categorization sits at the intersection of cognitive science and computer science: behavioral experiments examine how task structures shape human category representations~\cite{shepard1961learning,medin1978context,nosofsky1986attention,kruschke2020alcove}, while machine learning emphasizes maximizing classification accuracy through large-scale image-classification benchmarks such as ImageNet and deep architectures from AlexNet to ResNet~\cite{deng2009imagenet,russakovsky2015imagenet,krizhevsky2012imagenet,he2016deep}. To uncover the dynamics of category learning, recent work adopts three main complementary modeling frameworks: (1) \textbf{Reinforcement-learning (RL) models} treat learning as prediction-error–driven weight updates on an explicitly defined value function, capturing shifts in human rule weighting but relying on strong assumptions about states, values and update rules~\cite{sutton1998reinforcement,cohen2013high,cohen2021geometry}. (2) \textbf{Bayesian models} frame categorization as probabilistic inference over structured hypothesis spaces: hierarchical non-parametric priors (e.g., Chinese-restaurant or Dirichlet processes) compress the rule space, enabling one-shot generalization, latent-feature inference, and trial-by-trial predictions of feature-specific attention~\cite{sanborn2006more,tenenbaum2001generalization}. (3) Shallow \textbf{neural-network models} such as ALCOVE and SUSTAIN view category learning as an information processing stream--from feature input layer through attention weighting to response output layer--where attention weights progressively focus on diagnostic features~\cite{mack2020ventromedial, mack2016dynamic}. However, all these works reconstruct internal trajectories only retrospectively from behavior correlations, raising the question of whether such post-hoc reconstructions truly capture each individual’s internal states. This highlights the need for frameworks that can record and explain within-subject learning dynamics in real time~\cite{gershman2015computational,lieder2020resource}.